\documentclass[11pt]{article}

\usepackage[T1]{fontenc}
\usepackage[utf8]{inputenc}
\usepackage{lmodern}
\usepackage{microtype}
\usepackage{amsmath,amssymb,amsthm,mathtools}
\usepackage{aliascnt}
\usepackage{enumitem}
\usepackage{needspace}
\usepackage{longtable,array,booktabs}
\usepackage{xcolor}
\usepackage{xurl}
\usepackage{hyperref}
\usepackage[nameinlink,capitalise,noabbrev]{cleveref}

\emergencystretch=6em

\hypersetup{
  colorlinks=true,
  linkcolor=blue!45!black,
  citecolor=blue!45!black,
  urlcolor=blue!45!black
}

\newtheorem{theorem}{Theorem}[section]
\newaliascnt{lemma}{theorem}
\newtheorem{lemma}[lemma]{Lemma}
\aliascntresetthe{lemma}
\newaliascnt{proposition}{theorem}
\newtheorem{proposition}[proposition]{Proposition}
\aliascntresetthe{proposition}
\newaliascnt{corollary}{theorem}
\newtheorem{corollary}[corollary]{Corollary}
\aliascntresetthe{corollary}
\newaliascnt{definition}{theorem}
\newtheorem{definition}[definition]{Definition}
\aliascntresetthe{definition}
\newaliascnt{assumption}{theorem}
\newtheorem{assumption}[assumption]{Assumption}
\aliascntresetthe{assumption}
\newaliascnt{remark}{theorem}
\newtheorem{remark}[remark]{Remark}
\aliascntresetthe{remark}
\crefname{lemma}{Lemma}{Lemmas}
\crefname{proposition}{Proposition}{Propositions}
\crefname{corollary}{Corollary}{Corollaries}
\crefname{definition}{Definition}{Definitions}
\crefname{assumption}{Assumption}{Assumptions}
\crefname{remark}{Remark}{Remarks}

\newcommand{\HE}{Homotopical Encryption}
\newcommand{\HCP}{\mathrm{HCP}}
\newcommand{\CVP}{\mathrm{CVP}}
\newcommand{\MCA}{\mathrm{MCA}}
\newcommand{\RS}{\mathrm{RS}}
\newcommand{\F}{\mathbb{F}}
\newcommand{\R}{\mathbb{R}}
\newcommand{\Z}{\mathbb{Z}}
\newcommand{\Gr}{\operatorname{Gr}}
\newcommand{\Path}{\operatorname{Path}}
\newcommand{\Fib}{\operatorname{Fib}}
\newcommand{\image}{\operatorname{im}}
\newcommand{\coker}{\operatorname{coker}}
\newcommand{\dist}{\operatorname{dist}}
\newcommand{\Public}{\mathsf{Public}}
\newcommand{\Hidden}{\mathsf{Hidden}}
\newcommand{\Lean}{\textsc{Lean}}
\newcommand{\lean}[1]{\begingroup\Urlmuskip=0mu plus 1mu\path{#1}\endgroup}
\newcommand{\leanOmega}[1]{\ensuremath{\Omega}\lean{#1}}
\newenvironment{leanledger}{\small\raggedright}{\par}
\newcommand{\todo}[1]{\textcolor{red!65!black}{[TODO: #1]}}


\title{Homotopical Encryption: Hidden Hodge Operators and the Harmonic Cochain Problem}
\author{Chloe van der Vlugt}
\date{}

\hypersetup{
  pdftitle={Homotopical Encryption: Hidden Hodge Operators and the Harmonic Cochain Problem},
  pdfauthor={Chloe van der Vlugt}
}

\begin{document}

\maketitle

\begin{abstract}
We introduce the Harmonic Cochain Problem (HCP): recovering the operative
harmonic component \(h=P_{\ker(\Delta_1)}b\) from public data
\((\alpha,b)\). The degree-one Hodge--Laplacian \(\Delta_1\) is induced by
an operative hidden representative \(x\in A_\alpha=\pi^{-1}(\alpha)\) of
the public classifier value \(\alpha\). Its adjacent coboundaries
\(\delta_0,\delta_1\) determine the Hodge decomposition \(b=h+e\), with
\(e\in\image(\delta_0)\oplus\image(\delta_1^\ast)\).
We compare HCP with the closest-vector problem (CVP), distinguishing the
field-linear image \(\mathcal V_B(R)=\{Bx:x\in R^n\}\) from the discrete
lattice \(\Lambda_B=B\mathbb Z^n\). The former can be realized as a linear
kernel. For an explicitly supplied inclusion \(\Lambda_B\subseteq W\),
orthogonal projection onto the harmonic subspace \(W\) preserves the closest
lattice vectors. We then prove a hidden-operator separation: if publicly
equivalent hidden operators disagree on harmonicity, a correct known-lattice
CVP oracle cannot supply a uniformly exact public-data HCP solver.
\end{abstract}

\setcounter{tocdepth}{1}
\tableofcontents
\clearpage

\section{Introduction}

\subsection{Motivation and Contributions}

\HE{} studies recovery from public homotopy-class data when the Hodge operator
whose kernel is the message space is not part of the public input.  In a
known-lattice closest-vector problem, the lattice is presented before the
search begins.  In HCP, the recovery target is
\[
  h=P_H(b)\in\ker(\Delta_1),
\]
but the public class does not include the degree-one Hodge--Laplacian
\(\Delta_1\), its kernel, or the restriction maps from which \(\Delta_1\) is
built.  The private
recovery object is therefore not a scalar secret or a vector in a published
lattice; it is the harmonic cochain in the kernel of a hidden Hodge operator.

The operator dependence can be summarized as follows. A hidden representative
determines the restriction data, which determine the adjacent
coboundaries and hence the degree-one Hodge operator:
\[
  x
  \longmapsto
  \{R_{\sigma\to\tau}^{x}\}_{\sigma\prec\tau}
  \longmapsto
  (\delta_{0,x},\delta_{1,x})
  \longmapsto
  \Delta_{1,x}
  \longmapsto
  \bigl(\ker\Delta_{1,x},P_{H_x}\bigr).
\]
The recovery operation is Hodge-theoretic; ``homotopical'' describes the
framework in which multiple hidden representatives may occupy the same public
class while inducing different operators. The distinction is important: it is
hidden-operator ambiguity, not whether a displayed matrix happens to have
binary entries, that drives HCP.

This changes the comparison problem. A real linear span can be realized as the
kernel of a suitable linear operator, but a discrete lattice cannot equal a
nonzero real kernel. If an explicit lattice \(\Lambda\) is supplied inside a
harmonic subspace \(W=\ker(\Delta_W)\), orthogonal projection of the target
onto \(W\) preserves the closest vectors of \(\Lambda\). That bridge does not
make public-data HCP a known-lattice problem: public homotopy-class data
may be compatible with distinct hidden operators, distinct kernels, and
different harmonicity judgments.  The paper isolates this boundary and then
separates it from the finite verifier-facing projection artifacts used later.
The public classifier supplies the realized class value \(\alpha\).
The full cochain \(b=h+e\) is separate from the accepted exact commitment
payload: its decoded exact cochain has zero harmonic projection in its
specified Hodge system. A bound-observation relation ties the full sample to that payload,
its digest, snapshot and lifecycle context. Releasing \((\alpha,b)\)
requires an explicit disclosure or a specified extraction of that full
sample; a proof about the exact payload alone does not expose it.
The selected observation need not present \(\Delta_1\) or
\(\ker(\Delta_1)\). If full Pl\"ucker coordinates of the operative kernel
are additionally disclosed, however, they determine its fixed-metric
projector and hence the HCP target from \(b\).
The fibration ID is the spectral fingerprint
\(\operatorname{Spec}_{\mathrm{multi}}(\Delta_1)\), obtained only after the
relevant Hodge data have been recovered or certified; it is not part of the
public classifier value.

The principal notation is fixed globally as follows. Symbols used only inside
a single proposition or construction are defined locally.
\begin{table}[ht]
  \centering
  \small
  \begin{tabular}{@{}p{0.27\linewidth}p{0.65\linewidth}@{}}
    \hline
    Symbol & Meaning \\
    \hline
    \(p:E\to B\) & Helfer-style fibration projection \\
    \(\mathcal P_X(x,y)\) & path object from \(x\) to \(y\) in \(X\) \\
    \(\pi:A\to P\) & path-respecting public classifier \\
    \(\alpha=\pi(x)\in P\) & realized public classifier value \\
    \(A_\alpha=[x]_\pi\) & fiber of hidden representatives with public label \(\alpha\) \\
    \((\alpha,b)\) & public HCP observation \\
    \(\Delta_A:A\to A\times A\) & diagonal map in the fibrational semantics \\
    \(P_X\), \(F:P_X\to\mathsf{Vect}_{\R}\) & finite cell poset and its face-to-coface cellular sheaf \\
    \(a\in C^0,\ c=\delta_0a\in C^1\) & degree-zero opening and committed exact one-cochain \\
    \(\delta_0,\delta_1,\Delta_1\) & adjacent coboundaries and degree-one Hodge--Laplacian \\
    \(\ker(\Delta_1),\ P_H\) & harmonic subspace and harmonic projection \\
    \(\Delta_W=I-P_W\) & kernel-normalized projection operator associated with \(W\) \\
    \(\mathcal L,\ \ell\) & local syndrome target and affine MCA line, respectively \\
    \hline
  \end{tabular}
  \caption{Core notation used across the semantic and HCP layers.}
  \label{tab:core-notation}
\end{table}

This paper develops this boundary through the following statements.
\begin{enumerate}[label=(\roman*)]
  \item The Harmonic Cochain Problem (HCP): recovery of the operative
    component \(h=P_Hb\) for a hidden degree-one Hodge operator \(\Delta_1\).
  \item A Hodge-structured noise model with cochain decomposition
    \(b = h + e\), where \(h \in \ker(\Delta_1)\) and
    \(e \in \image(\delta_0) \oplus \image(\delta_1^\ast)\).
  \item A CVP comparison separating field-linear image from discrete lattice:
    generated subspaces can be realized as kernels, and, when
    \(\Lambda\subseteq W\) is explicitly supplied, projecting the target onto
    \(W\) preserves the discrete closest-vector condition on \(\Lambda\).
  \item A hidden-operator ambiguity theorem: publicly equivalent Hodge operators
    may disagree on harmonicity, so the CVP comparison does not make a correct
    explicit-lattice CVP oracle into a uniform public-data HCP solver.
  \item Constructed integral examples in the finite signed-incidence model:
    one homological public observation has distinct nonzero harmonic targets;
    a weighted pair shares \(b=(1,2)\) with integer coefficients two and one,
    with a unique CVP answer for each privately presented lattice.
  \item The supporting formal framework:
    fibrational identity, homotopy classes, exact Hodge/cochain invariants,
    spectral/Tsirelson admissibility, and Heyting-style public judgments.
  \item A projection chain through Pl\"ucker coordinates and cochain commitments to a
    canonical Reed--Solomon endpoint, with a conditional far-separation gate and
    an MCA bad-challenge bound imported at an explicit cited theorem boundary.
\end{enumerate}

Finite Hodge theory follows the cellular-cochain framework of
Hansen--Ghrist \cite{hansenGhrist2018sheaves}; the algebraic model used here is
specified in \cref{sec:hcp-definition}. Kernel realization and the
Pythagorean projection identity are elementary linear-algebraic foundations.
Helfer's homotopy-invariant semantics \cite{helfer2019fol} motivates the public
classifier; the constructed examples prove only the explicitly named
nonzero-cohomology-test fragment. The selected formalization integrates these
objects with the constructed observation ambiguities and recovery boundaries.
It establishes neither an asymptotic CVP-hardness reduction nor a sampled
encryption-security theorem. The independently stated Lean entry covers the
six Hodge/CVP groups and two observation-boundary groups in its source
crosswalk; the broader spectral, oracle and coding interfaces retain their
separately stated assumptions and formalization status.

The companion manuscript \emph{Capacity-Radius Proximity for Smooth-Domain
Reed--Solomon Codes} \cite{vanDerVlugt2026proximity} develops the universal
list-decoding and MCA results independently of HCP. The present paper supplies
the HCP foundations and the application-specific projection interface. Shared
coding definitions connect the two papers; the companion's coding theorems do
not supply the source-to-code far-separation witness required here.

The corresponding Lean definitions, reductions, gates, theorem boundaries,
paper-to-Lean cross-reference, and file-by-file source coverage are recorded in
the appendices.

\section{Primitives}\label{sec:primitives}

Let \(A\) be the type of hidden presentations, with semantic interpretation
\(M_x\) for each \(x\in A\). Write \(x\simeq y\) for homotopy equivalence
of \(M_x\) and \(M_y\) in Helfer's sense \cite[Definition~7.1]{helfer2019fol}.
Assume that this equivalence preserves the chosen closed tests. Let \(P\) be
the type of public class data and choose a classifier \(\pi:A\to P\) that factors through their
closed theory. Then
\[
  x\simeq y\quad\Longrightarrow\quad\pi(x)=\pi(y).
\]
The typed data and the invariance premise are made precise in
\cref{def:typed-hcp-realization,prop:helfer-hcp-bridge}.
For a realized public value \(\alpha\), write
\[
  \alpha\in\image(\pi),\qquad
  A_\alpha=\pi^{-1}(\alpha)=\{x\in A:\pi(x)=\alpha\}.
\]
For \(x\in A\), we also write
\([x]_\pi:=A_{\pi(x)}\). Thus, whenever \(\alpha=\pi(x)\), the two
descriptions agree: \(A_\alpha=[x]_\pi\). We call the classifier
\emph{complete for the declared homotopy equivalence} when
\(\pi(x)=\pi(y)\) implies \(x\simeq y\) for all \(x,y\in A\).

We attach Hodge data to each presentation using finite graded cell carriers,
integer incidence numbers and face-to-coface linear stalk maps.

Fix a finite coordinate set \(N\) and let \(V=\mathbb R^N\) have its Euclidean
inner product. Each presentation has finite-dimensional real inner-product
cochain spaces \(C_x^0,C_x^1,C_x^2\). To compare presentations using one public
cochain vector, fix isometric coordinate charts
\(\iota_x:C_x^1\overset{\sim}{\longrightarrow}V\).
We identify degree-one vectors and maps with their transported versions in
\(V\); the charts are part of the comparison data.

The restriction maps \(R^x_{\sigma\to\tau}\) for \(\sigma\prec\tau\),
including vertex-to-edge and edge-to-face maps, determine adjacent
coboundaries
\[
  \delta_{0,x}:C_x^0\to V,
  \qquad
  \delta_{1,x}:V\to C_x^2,
  \qquad
  \delta_{1,x}\delta_{0,x}=0.
\]
In these coordinates, define the degree-one Hodge--Laplacian, its harmonic
subspace and its orthogonal projector by
\[
\begin{aligned}
  \Delta_x&=\delta_{0,x}\delta_{0,x}^\ast+
    \delta_{1,x}^\ast\delta_{1,x},\\
  \mathcal H_x&=\ker(\Delta_x),\qquad
    P_{\mathcal H_x}:V\to\mathcal H_x\subseteq V.
\end{aligned}
\]
The stars denote adjoints. When the operative representative is fixed, we
suppress \(x\) and also write \(\Delta_1\) or \(\Delta\) for its degree-one
Laplacian, and \(P_H\) for its harmonic projector.

The algebraic cellular model above adapts the finite Hodge setup of
\cite{hansenGhrist2018sheaves} without assuming a regular CW or simplicial
realization for every admitted incidence structure. Such a realization and
the separate simplicial-coherence conditions require additional data.

\begin{definition}[Matrix harmonicity]\label{def:matrix-harmonicity}
  For a hidden linear Hodge operator \(\Delta\), a cochain
  \(h\in\mathbb R^N\)
  is \(\Delta\)-harmonic when
  \[
    \Delta h=0.
  \]
  Thus the harmonic kernel is \(\ker(\Delta)\). In the projection model this
  agrees with the condition \(h\in W\) for the kernel-normalized operator
  \(\Delta_W=I-P_W\).
\end{definition}

\begin{definition}[Harmonic projection]\label{def:projection-laplacian}
  For a subspace \(W\le V\), let \(P_W:V\to W\subset V\) denote the
  orthogonal projection. Its associated kernel-normalized operator is
  \(\Delta_W=I-P_W\), with \(\ker\Delta_W=W\).
\end{definition}

\begin{proposition}[Hodge decomposition surface]\label{prop:hodge-decomposition}
  Let the three consecutive cochain spaces be finite-dimensional real inner
  product spaces, and let the displayed stars denote adjoints. Suppose
  \[
    d_{\mathrm{next}}d_{\mathrm{prev}}=0,
    \qquad
    \Delta=d_{\mathrm{prev}}d_{\mathrm{prev}}^\ast
      +d_{\mathrm{next}}^\ast d_{\mathrm{next}} .
  \]
  Then every cochain \(b\) admits a decomposition
  \[
    b=h+e_1+e_2
  \]
  with
  \[
    \Delta h=0,\qquad
    e_1\in\image(d_{\mathrm{prev}}),\qquad
    e_2\in\image(d_{\mathrm{next}}^\ast),
  \]
  and
  \[
    \langle h,e_1\rangle=\langle h,e_2\rangle=\langle e_1,e_2\rangle=0.
  \]
  The decomposition is unique.
\end{proposition}

\begin{proof}
  Each summand defining \(\Delta\) is self-adjoint, so \(\Delta\) is
  self-adjoint. In finite dimension,
  \(\image(\Delta)=(\ker\Delta)^\perp\), and hence the ambient cochain
  space decomposes as
  \[
    \ker(\Delta)\oplus\image(\Delta).
  \]
  Since
  \[
    \image(\Delta)
    \subseteq
    \image(d_{\mathrm{prev}})+\image(d_{\mathrm{next}}^\ast),
  \]
  every cochain \(b\) can be written as \(h+e_1+e_2\) with \(h\in\ker\Delta\),
  \(e_1\in\image(d_{\mathrm{prev}})\), and
  \(e_2\in\image(d_{\mathrm{next}}^\ast)\). From \(\Delta h=0\), one obtains
  \[
    \|d_{\mathrm{prev}}^\ast h\|^2+\|d_{\mathrm{next}}h\|^2=0,
  \]
  hence \(d_{\mathrm{prev}}^\ast h=0\) and \(d_{\mathrm{next}}h=0\). These
  identities give \(\langle h,e_1\rangle=0\) and
  \(\langle h,e_2\rangle=0\), while \(d_{\mathrm{next}}d_{\mathrm{prev}}=0\)
  gives \(\langle e_1,e_2\rangle=0\). Hence the harmonic, exact, and
  coexact subspaces are pairwise orthogonal. Their sum is therefore direct,
  which makes the three components unique.
\end{proof}

\subsection{Hodge Encryption and Correctness}

For a hidden representative \(x\), define its noise space by
\[
  \mathcal E_x=\mathcal H_x^\perp.
\]
The Hodge decomposition identifies the noise space as the orthogonal direct
sum
\[
  \mathcal E_x
  =
  \image(\delta_{0,x})
  \oplus^\perp
  \image(\delta_{1,x}^{\ast}),
\]
and gives the orthogonal decomposition
\[
  C_x^1=\mathcal H_x\oplus^\perp\mathcal E_x.
\]
Consequently,
\[
\begin{aligned}
  \operatorname{combine}_x:
    \mathcal H_x\times\mathcal E_x&\longrightarrow C_x^1,
  &(h,e)&\longmapsto h+e,\\
  \operatorname{split}_x:C_x^1&\longrightarrow
    \mathcal H_x\times\mathcal E_x,
  &b&\longmapsto
    \bigl(P_{\mathcal H_x}b,P_{\mathcal E_x}b\bigr)
\end{aligned}
\]
are inverse linear bijections. Orthogonality also gives
\[
  \|h+e\|^2=\|h\|^2+\|e\|^2.
\]

This bijection is the algebraic encryption/decryption construction. Given a
harmonic message \(h\in\mathcal H_x\), randomized Hodge encryption samples, or
otherwise chooses, admissible noise \(e\in\mathcal E_x\) and sets
\[
  \mathsf{Enc}_x(h;e)=
  \operatorname{combine}_x(h,e)=h+e.
\]
Full coordinate recovery is \(\operatorname{split}_x\); message decryption
retains its harmonic coordinate:
\[
  \mathsf{Dec}_x(b)
  =
  \operatorname{pr}_{\mathcal H_x}
    \bigl(\operatorname{split}_x(b)\bigr)
  =
  P_{\mathcal H_x}b.
\]
For every admissible noise value,
\[
  \mathsf{Dec}_x\bigl(\mathsf{Enc}_x(h;e)\bigr)=h,
  \qquad
  \operatorname{split}_x\bigl(\operatorname{combine}_x(h,e)\bigr)=(h,e).
\]
The bijection concerns the full message--noise pair \((h,e)\). Choosing a
noise distribution makes \(\mathsf{Enc}_x\) randomized, while
\(\mathsf{Dec}_x\) recovers the message for every admissible noise value.
Both maps depend on the operative representative \(x\).

In the \HE{} public-sample model, a presentation \(x\) produces a cochain
\[
  b=h+e,
  \qquad
  h\in\mathcal H_x,
  \qquad
  e\in\image(\delta_{0,x})\oplus\image(\delta_{1,x}^\ast).
\]
The message \(h\) is a harmonic representative of its degree-one cohomology
class, and \(e\) is Hodge-structured noise. When \(x\) is also equipped with
an explicit harmonic-lattice presentation, write
\(\Lambda_x\subseteq\mathcal H_x\) for its discrete lattice and choose
\(h\in\Lambda_x\). In this integral profile, the public fiber indexes a
family of possible embedded harmonic lattices; the operative hidden member
selects one of them.

\section{The Harmonic Cochain Problem}\label{sec:hcp-definition}

\begin{definition}[Public observation]\label{def:public-observation}
  A public HCP observation is a pair \((\alpha,b)\) consisting of a realized
  public classifier value \(\alpha\in\image(\pi)\) and an observed degree-one
  cochain \(b\in V\) in the fixed comparison coordinates.
\end{definition}

\begin{definition}[Hidden Hodge witness]\label{def:hidden-hodge-witness}
  A hidden Hodge witness for a public observation \((\alpha,b)\) is a hidden
  representative \(x\in A_\alpha\), satisfying \(\pi(x)=\alpha\). Its
  presentation data select the operative operator \(\Delta_x\). The witness
  and \(b\) determine the coordinates
  \[
    h_x(b):=P_{\mathcal H_x}b,
    \qquad
    e_x(b):=b-h_x(b).
  \]
  The Hodge decomposition in \cref{prop:hodge-decomposition} gives
  \[
    h_x(b)\in\ker(\Delta_x),\qquad
    e_x(b)\in
      \image(\delta_{0,x})\oplus\image(\delta_{1,x}^\ast),\qquad
    b=h_x(b)+e_x(b).
  \]
\end{definition}

\begin{definition}[Harmonic Cochain Problem]\label{def:hcp}
  The \emph{Harmonic Cochain Problem} (HCP) asks, given the public observation
  \((\alpha,b)\), to recover the operative harmonic component
  \[
    h=h_x(b)\in\mathcal H_x,
  \]
  equivalently the harmonic term in the orthogonal Hodge decomposition of
  \(b\) for the operative hidden witness \(x\in A_\alpha\). The witness is
  bound to the instance but is not part of the public input.
\end{definition}

\begin{definition}[HCP solution]\label{def:hcp-solution}
  Relative to an operative hidden witness \(x\), a solution to an
  HCP instance \((\alpha,b)\) is the cochain
  \(h'=h_x(b)=P_{\ker(\Delta_x)}b\). A public-data-only solver is
  \emph{uniformly exact}
  if the same output has this value for every hidden witness compatible with
  its public input.
\end{definition}

\subsection{Public Classes and Hidden Operators}\label{sec:mathematical-framework}

\begin{definition}[Hidden-operator public projection]\label{def:hidden-operator-public-projection}
  For a public class type \(P\), a public projection is a map
  \[
    \pi_\Delta:\{\text{hidden Hodge operators }\Delta\}\longrightarrow P.
  \]
  Two hidden operators are publicly equivalent when
  \(\pi_\Delta(\Delta)=\pi_\Delta(\Delta')\). The hidden operator fiber over
  \(\alpha\in P\) is
  \[
    \operatorname{Fib}_{\pi_\Delta}(\alpha)
    =\pi_\Delta^{-1}(\alpha)
    =\{\Delta:\pi_\Delta(\Delta)=\alpha\},
  \]
  and the hidden class of \(\Delta\) is
  \(\operatorname{Fib}_{\pi_\Delta}(\pi_\Delta(\Delta))\).
\end{definition}

\begin{definition}[Lossiness and kernel ambiguity]\label{def:kernel-ambiguity}
  The public projection \(\pi_\Delta\) is \emph{lossy} if there exist
  \(\Delta\ne\Delta'\) with
  \(\pi_\Delta(\Delta)=\pi_\Delta(\Delta')\). The two
  operators have the same harmonic kernel if
  \[
    \forall v,\quad \Delta v=0\Longleftrightarrow \Delta'v=0.
  \]
  The projection is \emph{kernel-ambiguous} if there are publicly equivalent
  operators \(\Delta,\Delta'\) whose harmonic kernels are not the same.
\end{definition}

\begin{definition}[Public-compatible Hodge ambiguity]\label{def:public-compatible-hodge-ambiguity}
  Public-compatible Hodge ambiguity means that there exist publicly equivalent
  operators \(\Delta,\Delta'\) and a cochain-facing observation \(b\) such that
  \[
    \Delta b=0
    \qquad\text{but}\qquad
    \Delta'b\ne 0.
  \]
\end{definition}

\begin{definition}[Typed HCP semantic realization]\label{def:typed-hcp-realization}
  Fix a multi-sorted algebraic signature \(\sigma\), a model category
  \(\mathcal C\) with terminal object and binary products, and, for each
  hidden representative \(x\in A\), a base
  interpretation \(M_x:\sigma\to\mathcal C_{\mathrm f}\) together with
  homotopical formula semantics over \(M_x\). Write
  \[
    \operatorname{Th}(M_x)
    =\{\varphi\text{ closed}:M_x\vDash\varphi\}
  \]
  for its closed first-order theory. Relative to the presentation data and
  coordinate charts fixed above, a typed HCP semantic realization consists
  additionally of maps
  \[
  \begin{aligned}
    x&\longmapsto\Delta_x\in\operatorname{End}(V),\\
    \kappa&:\{\text{closed theories}\}\longrightarrow P,\\
    \pi_\Delta&:\{\text{Hodge operators}\}\longrightarrow P,
  \end{aligned}
  \]
  satisfying the two factorization identities
  \[
    \pi(x)=\kappa(\operatorname{Th}(M_x)),
    \qquad
    \pi_\Delta(\Delta_x)=\pi(x).
  \]
  The Laplacian \(\Delta_x\) is presentation-level Hodge data extending the
  semantic core \(M_x\); it is not silently included among the symbols whose
  closed theory is asserted to be homotopy invariant.
\end{definition}

\begin{proposition}[Helfer-to-HCP classifier bridge]\label{prop:helfer-hcp-bridge}
  Let \(x,y\in A\) belong to a typed HCP semantic realization. If
  \(M_x\) and \(M_y\) are homotopy equivalent interpretations in the sense of
  Helfer, then
  \[
    \pi(x)=\pi(y)
    \qquad\text{and}\qquad
    \pi_\Delta(\Delta_x)=\pi_\Delta(\Delta_y).
  \]
  If, independently, there is a cochain \(b\) for which
  \[
    \Delta_xb=0,
    \qquad
    \Delta_yb\ne0,
  \]
  then \(\pi_\Delta\) has public-compatible Hodge ambiguity, and hence is
  kernel-ambiguous and lossy.
\end{proposition}

\begin{proof}
  The stated closed-test invariance premise gives, for every closed
  \(\sigma\)-formula \(\varphi\),
  \[
    M_x\vDash\varphi
    \quad\Longleftrightarrow\quad
    M_y\vDash\varphi.
  \]
  In the full first-order setting this premise is supplied by Helfer's
  homotopy-invariance theorem \cite[Theorem~12.3]{helfer2019fol}.
  Hence \(\operatorname{Th}(M_x)=\operatorname{Th}(M_y)\). Applying
  \(\kappa\) and then the two factorization identities gives
  \[
    \pi(x)=\pi(y),
    \qquad
    \pi_\Delta(\Delta_x)=\pi_\Delta(\Delta_y).
  \]
  The additional displayed harmonicity and non-harmonicity conditions now
  exhibit \((\Delta_x,\Delta_y,b)\) as a public-compatible Hodge-ambiguity
  witness. The ambiguity hierarchy below supplies kernel ambiguity and
  lossiness.
\end{proof}

Under closed-test invariance and classifier completeness,
\cref{prop:helfer-hcp-bridge} gives
\[
  \pi(x)=\pi(y)\quad\Longleftrightarrow\quad x\simeq y,
  \qquad
  A_{\pi(x)}=\{y\in A:y\simeq x\}.
\]
The implication from label equality uses completeness; the other uses
invariance. The fiber identity follows from the definition of \(A_{\pi(x)}\).
The family in \cref{prop:cellular-axis-ambiguity} has a complete classifier.
Completeness concerns the semantic class; it does not select its operative
embedded harmonic kernel. Kernel disagreement remains an independent condition. If
\(\Delta\)-harmonicity were included in the invariant first-order signature,
and the common coordinate \(b\) were
named compatibly in both interpretations, Helfer's theorem would instead rule
out the displayed disagreement.

Suppose an equivalence of hidden presentations induces degreewise linear
isometric equivalences
\(U_0,U_1,U_2\) satisfying
\[
  U_1\delta_{0,x}=\delta_{0,y}U_0,
  \qquad
  U_2\delta_{1,x}=\delta_{1,y}U_1.
\]
In the fixed coordinate charts, \(U_1\) is an isometry of \(V\). The
isometries transport adjoints, giving
\[
  \Delta_yU_1=U_1\Delta_x,
  \qquad
  U_1(\ker\Delta_x)=\ker\Delta_y.
\]
The corresponding orthogonal projections commute with transport:
\[
  U_1P_{\ker\Delta_x}b
  =
  P_{\ker\Delta_y}U_1b.
\]
Thus a transported observation preserves both the operative target and
harmonicity:
\[
  \Delta_xb=0
  \quad\Longleftrightarrow\quad
  \Delta_y(U_1b)=0.
\]
The ambiguity witness in \cref{prop:helfer-hcp-bridge} uses a fixed public
coordinate \(b\in V\). Under the displayed conjugacy, its separation condition is
equivalently
\[
  b\in\ker\Delta_x,
  \qquad
  b\notin U_1(\ker\Delta_x).
\]
It requires a nonstabilization witness in addition to the equivalence, with
the private transport \(U_1\) neither revealed nor applied.

\begin{proposition}[Ambiguity hierarchy]\label{prop:ambiguity-hierarchy}
  Public-compatible Hodge ambiguity implies kernel ambiguity, and kernel
  ambiguity implies that the public projection is lossy.
\end{proposition}

\begin{proof}
  If \(b\) is harmonic for \(\Delta\) but not for \(\Delta'\), then the two
  operators cannot have the same harmonic kernel. This gives kernel ambiguity.
  If publicly equivalent operators have different kernels, they cannot be the
  same operator, since equality of operators implies equality of kernels.
  Therefore the public projection identifies two distinct hidden operators.
\end{proof}

\subsection{Computational Scope}\label{sec:computational-scope}

These exact-recovery comparisons do not by themselves give an asymptotic
CVP-hardness reduction. Such a comparison requires encoded instance families,
an input-size measure, an algorithm class and a success criterion. A
distributional security claim also requires a sampling model and a declared
adversary view.

\section{CVP Comparison and Kernel Ambiguity}

CVP asks for a closest vector in an explicitly presented lattice
\cite{micciancioGoldwasser2002}. HCP asks for
the projection of \(b\) into the kernel of the operative hidden Hodge operator.

\subsection{Known-Lattice CVP}

An explicit Euclidean CVP instance consists of a full-column-rank generator
matrix \(B\in\mathbb R^{m\times n}\) and a target \(t\in\mathbb R^m\). Its
discrete lattice is
\[
  \Lambda_B=B\mathbb Z^n=\{Bz:z\in\mathbb Z^n\}.
\]
The injectivity of \(z\mapsto Bz\) records that the displayed columns are a
lattice basis. A correct closest-vector answer is a vector
\(v\in\Lambda_B\) such that
\[
  \|t-v\|^2\le \|t-w\|^2
  \qquad
  \text{for every }w\in\Lambda_B.
\]
Thus, for fixed dimensions, a CVP oracle \(O\) is correct when
\[
\begin{aligned}
  \mathsf{CVP}_{\mathrm{corr}}(O)
  &\quad\Longleftrightarrow\quad
  \forall\text{ full-column-rank }B,t,\\
  &\qquad O(B,t)\in\Lambda_B
  \quad\text{and}\quad
  \forall w\in\Lambda_B,\\
  &\qquad \|t-O(B,t)\|^2\le \|t-w\|^2.
\end{aligned}
\]
The weaker output condition is merely lattice membership:
\[
  \forall B,t,\ O(B,t)\in\Lambda_B .
\]
Fix a class \(\mathcal C\) of admissible algorithms, such as polynomial-time or
probabilistic polynomial-time algorithms. The external CVP hardness assumption
is the schema
\[
\begin{aligned}
  \mathsf{CVP}_{\mathrm{hard}}(\mathcal C)
  \quad\Longleftrightarrow\quad
  \neg\exists O,\ 
    O\in\mathcal C\wedge\mathsf{CVP}_{\mathrm{corr}}(O).
\end{aligned}
\]

\subsection{Presented Linear Images}

Let \(R\) be a field. The field-linear image of \(B\) is a different object:
\[
  \mathcal V_B(R)=BR^n=\{Bx:x\in R^n\}.
\]
Over a field this is a linear subspace, not the discrete lattice
\(B\mathbb Z^n\).

Let \(B\) present \(\mathcal V_B(R)\). If a matrix
\(\Delta\) satisfies \(\ker(\Delta)=\mathcal V_B(R)\), then any solver satisfying
\[
  \Delta S(\Delta,t)=0
\]
returns \(S(\Delta,t)\in\mathcal V_B(R)\).

\begin{proposition}[Kernel realization for explicit generated subspaces]\label{prop:hcp-to-cvp}
  Let \(B\in R^{n\times n}\) be an explicit square generator matrix over a
  field and set
  \(\mathcal V_B(R)=\{Bx:x\in R^n\}\). There exists a matrix \(\Delta\) such that
  \[
    \ker(\Delta)=\mathcal V_B(R) .
  \]
  Consequently, any solver \(S(\Delta,b)\) whose output is
  \(\Delta\)-harmonic for every
  presented \(\Delta\) returns a vector in the explicit generated subspace
  \(\mathcal V_B(R)\) when run on this realized operator.
\end{proposition}

\begin{proof}
  Let \(B\in R^{n\times n}\) be an explicit generator matrix over a field
  \(R\), set
    \[
      \mathcal V_B(R)=\{Bx:x\in R^n\}\subseteq R^n,
    \]
    and let \(t\in R^n\) be the solver input. We construct a matrix \(\Delta\)
    with
  \[
    \ker(\Delta)=\mathcal V_B(R).
  \]
  Put \(V=\mathcal V_B(R)\). Since \(R^n\) is finite-dimensional over a field,
  choose a complementary subspace \(W\) with
  \[
    V+W=R^n,
    \qquad
    V\cap W=0 .
  \]
  Thus every \(y\in R^n\) has a unique decomposition
  \[
    y=v+w,
    \qquad
    v\in V,\quad w\in W.
  \]
  Equivalently, there is a unique element \(w_y\in W\) such that
  \(y-w_y\in V\). Define
  \[
    P(y)=w_y.
  \]
  The uniqueness of the \(V\oplus W\) decomposition gives linearity: for
  \(y,z\in R^n\), the vector \(P(y)+P(z)\) lies in \(W\), and
  \[
    (y+z)-(P(y)+P(z))
    =
    (y-P(y))+(z-P(z))\in V,
  \]
  so \(P(y+z)=P(y)+P(z)\). Similarly, for \(a\in R\),
  \(aP(y)\in W\) and
  \[
    ay-aP(y)=a(y-P(y))\in V,
  \]
  so \(P(ay)=aP(y)\).

  We now compute the kernel. If \(y\in V\), then \(0\in W\) and \(y-0\in V\),
  hence uniqueness gives \(P(y)=0\). Conversely, if \(P(y)=0\), then
  \(y-P(y)=y\in V\). Therefore
  \[
    \ker P=V=\mathcal V_B(R).
  \]
  Let \(\Delta\) be the matrix of the linear map \(P\). Then, for every
  \(y\in R^n\),
  \[
    \Delta y=0
    \quad\Longleftrightarrow\quad
    y\in\ker P
    \quad\Longleftrightarrow\quad
    y\in\mathcal V_B(R)
    \quad\Longleftrightarrow\quad
    \exists x,\ Bx=y .
  \]

    Let \(S:R^{n\times n}\times R^n\to R^n\) be a presented-operator Hodge
  solver satisfying
  \[
    \forall \Delta\, b,\qquad
    \Delta(S(\Delta,b))=0.
  \]
    Applying this to the \(\Delta\) constructed above and \(b=t\) gives
  \[
    \Delta(S(\Delta,t))=0.
  \]
  By the kernel calculation, there exists \(x\) with
  \[
    Bx=S(\Delta,t).
  \]
    Hence \(S(\Delta,t)\in\mathcal V_B(R)\).
\end{proof}

When \(R=\mathbb R\), if a presented-operator solver additionally minimizes
Euclidean distance among \(\Delta\)-harmonic vectors, the same kernel equality
says that it returns the continuous closest point of \(\mathcal V_B(\mathbb
R)\). This is still not discrete CVP unless a separate lattice is supplied.

\subsection{Explicit Harmonic-Lattice Bridge}

The discrete bridge requires both kinds of data: a real harmonic subspace and
an integer lattice contained in it. Let \(W\subseteq\mathbb R^m\) be the
operative harmonic subspace, let \(P_W\) be its orthogonal projection, and let
\(\Lambda_B=B\mathbb Z^n\subseteq W\) be an explicitly presented lattice. Set
\[
  \Delta_W=I-P_W.
\]
The inclusion \(\Lambda_B\subseteq W\) does not identify the discrete lattice
with the real subspace.

\begin{proposition}[Discrete CVP through harmonic projection]\label{prop:hodge-projection-is-cvp}
  For \(t\in\mathbb R^m\), put \(t_W=P_Wt\). Then:
  \begin{enumerate}[label=(\roman*)]
    \item \(t_W\) and every \(\lambda\in\Lambda_B\) are
      \(\Delta_W\)-harmonic;
    \item for every \(\lambda\in\Lambda_B\),
      \[
        \|t-\lambda\|^2
        =
        \|t-t_W\|^2+\|t_W-\lambda\|^2;
      \]
    \item a lattice vector \(\lambda\) is closest to \(t\) in
      \(\Lambda_B\) if and only if it is closest to \(t_W\) in
      \(\Lambda_B\).
  \end{enumerate}
\end{proposition}

\begin{proof}
  Since \(t_W\in W\), every \(\lambda\in\Lambda_B\) also lies in \(W\), and
  \(\ker(\Delta_W)=W\), the harmonicity claims follow. For the distance
  identity, write
  \[
    t-\lambda=(t-t_W)+(t_W-\lambda).
  \]
  Orthogonal projection gives \(t-t_W\in W^\perp\), while
  \(t_W-\lambda\in W\). The two summands are orthogonal, so the Pythagorean
  identity gives (ii). The term \(\|t-t_W\|^2\) is independent of
  \(\lambda\); consequently, for any \(\lambda,\mu\in\Lambda_B\),
  \[
    \|t-\lambda\|^2\le \|t-\mu\|^2
    \quad\Longleftrightarrow\quad
    \|t_W-\lambda\|^2\le \|t_W-\mu\|^2.
  \]
  This is exactly (iii).
\end{proof}

The continuous projection \(t_W=P_Wt\) is the closest harmonic vector to
\(t\), but it need not lie in \(\Lambda_B\). Thus projection transports the
target to the harmonic subspace while preserving the discrete minimizers; it
does not perform the remaining lattice search.

\subsection{Oracle Insufficiency}

Fix a typed HCP semantic realization as in \cref{def:typed-hcp-realization}.
The solver contracts below range over its admitted representatives.

\begin{definition}[Public-data Hodge solver contract]\label{def:public-hodge-solver-contract}
  Let
  \[
    S:\mathcal O_{\mathrm{CVP}}\times
      (P\times\mathbb R^n)\longrightarrow\mathbb R^n,
  \]
  where \(\mathcal O_{\mathrm{CVP}}\) is the class of explicit-lattice CVP
  oracles. For a realization \(x\mapsto(M_x,\Delta_x)\) and
  \(O\in\mathcal O_{\mathrm{CVP}}\), define
  \[
  \begin{aligned}
    \mathsf H_O(S)
    &\Longleftrightarrow
    \forall x:A,\quad
    \forall b\in\mathbb R^n,\quad
    \Delta_x\bigl(S(O,(\pi(x),b))\bigr)=0,\\
    \mathsf F_O(S)
    &\Longleftrightarrow
    \forall x:A,\quad
    \forall b\in\mathbb R^n,\quad
    \Delta_xb=0\Rightarrow S(O,(\pi(x),b))=b.
  \end{aligned}
  \]
  The public-data-only contract is
  \[
    \exists S,\quad
    \forall O,\quad
    \mathsf{CVP}_{\mathrm{corr}}(O)\Rightarrow
    \bigl(
      \mathsf H_O(S)
      \wedge
      \mathsf F_O(S)
    \bigr).
  \]
\end{definition}

\begin{proposition}[Typed HCP ambiguity and CVP oracle insufficiency]\label{prop:cvp-insufficient}
  Let \(x,y\in A\) be representatives in a typed HCP semantic realization.
  Suppose that \(M_x\) and \(M_y\) are homotopy equivalent and that some
  \(b\in\mathbb R^n\) satisfies
  \[
    \Delta_xb=0,
    \qquad
    \Delta_yb\ne0.
  \]
  If at least one correct explicit-lattice CVP oracle exists, then
  \[
    \neg\exists S,\quad
    \forall O,\quad
    \mathsf{CVP}_{\mathrm{corr}}(O)\Rightarrow
    \bigl(
      \mathsf H_O(S)
      \wedge
      \mathsf F_O(S)
    \bigr).
  \]
  Hence a correct oracle for explicit known-lattice CVP does not by itself
  determine the operative Hodge--Laplacian kernel.
\end{proposition}

\begin{proof}
  By \cref{prop:helfer-hcp-bridge},
  \(\pi(x)=\pi(y)\). Suppose
  \(\exists O,\mathsf{CVP}_{\mathrm{corr}}(O)\), and choose such an \(O\).
  Suppose, for contradiction, that \(S\) satisfies the displayed
  public-data-only contract. Then
  \[
    \mathsf H_O(S)\wedge\mathsf F_O(S).
  \]
  Since \(\Delta_xb=0\), exactness at \(x\) gives
  \[
    S(O,(\pi(x),b))=b.
  \]
  Soundness at \(y\) gives
  \[
    \Delta_yS(O,(\pi(y),b))=0.
  \]
  Replacing \(\pi(y)\) by \(\pi(x)\) and using the exactness equality yields
  \(\Delta_yb=0\), contradicting the separation hypothesis.
\end{proof}

For every dimension at least two, consider the elementary matrix witness \(L=0\),
\(L'=\operatorname{diag}(1,0,\ldots,0)\), and \(b=e_1\) under the constant
projection to a one-point public type. This is a sanity example showing that
the abstract ambiguity predicates are inhabited. It is not used as evidence
that an arbitrary HCP realization admits the required semantic and Hodge
extension witness. That witness remains an explicit hypothesis of
\cref{prop:cvp-insufficient}; the following family constructs one for its
stated semantic fragment.

\subsection{A Constructed Integral Cellular Family}

The preceding implication has a concrete instantiation in a stated
homological fragment of the typed Helfer interface. Take one vertex, one edge
and one face with incidence numbers \(\partial_1=0\), \(\partial_2=1\),
and stalks \(0,\mathbb R^2,\mathbb R\), respectively. This is the algebraic
finite signed-incidence model of \cref{sec:primitives}. Two cellular
sheaves in that model have
\[
  \delta_{0,0}=\delta_{0,1}=0,
  \qquad
  \delta_{1,0}(x_0,x_1)=x_0,
  \qquad
  \delta_{1,1}(x_0,x_1)=x_1.
\]
Their true degree-one Hodge operators are
\[
  \Delta_0=\operatorname{diag}(1,0),
  \qquad
  \Delta_1=\operatorname{diag}(0,1).
\]
Coordinate exchange is a degreewise chain isometry. Both operators have
positive spectral gap \(1\). Standard integral stalk lattices give harmonic
lattices \(\mathbb Z e_1\) and \(\mathbb Z e_0\), each with an explicit
full-span presentation.

Interpret these degree-zero-through-two complexes in the bounded-below
real cochain Quillen model category, with the incidence-defined coboundaries
displayed above. The closed tests are indexed by integers and ask whether
cohomology in that degree is nonzero. The required cofibrancy and fibrancy
hold for these real complexes. The model's weak equivalences are
quasi-isomorphisms, which preserve the tests and give the required
closed-test invariance under the declared typed Helfer equivalences.

The coordinate-exchange chain isometry gives the same homological theory
to both presentations. On this admitted family, equality of the public
theories is equivalent to typed homotopy equivalence.

\Needspace{10\baselineskip}
\begin{proposition}[One observation, two nonzero integral harmonic targets]
  \label{prop:cellular-axis-ambiguity}
  Let \(p\) be the common closed theory of these two presentations in the
  stated fragment, and let \(b=e_0+e_1\). There are structured Hodge samples
  with that same observation \((p,b)\), harmonic messages \(e_1\) and \(e_0\),
  and zero exact components. Their operative targets are distinct and nonzero.
  For each presentation's explicit harmonic lattice, its message is also the
  unique closest lattice vector to \(b\).
\end{proposition}

\begin{proof}
  In the first presentation use harmonic message \(e_1\), exact opening zero
  and coexact opening \(1\), producing noise \(e_0\). In the second use
  message \(e_0\), exact opening zero and coexact opening \(1\), producing
  noise \(e_1\). Both sums equal \(b\), and both messages belong to their
  standard integral harmonic lattices. The actual Hodge projections recover
  those messages. Orthogonality gives the Pythagorean identity from
  \cref{prop:hodge-projection-is-cvp}; any other lattice point has strictly
  positive squared distance from the recovered message. Thus each message
  is the unique CVP answer in its presented lattice.
\end{proof}

For the same realization, the fixed coordinate \(e_1\) is harmonic for
\(\Delta_0\) and not for \(\Delta_1\). This constructs the separation witness
used in \cref{prop:cvp-insufficient}. Equality of the public theory is
derived from the
actual semantic equivalence, rather than imposed by a one-point classifier.
The matrix-level public projection is extended from the two admitted
operators; no semantic realization is asserted for other matrices.

The displayed pair uses coefficient \(1\) in each presentation's own
harmonic-lattice basis. Its distinct ambient targets therefore do not alone
assert distinct encoded integer plaintexts. A second constructed pair makes
the encoding explicit: take \(b=e_0+2e_1\), harmonic messages \(2e_1\) and
\(e_0\), and complementary coexact noise \(e_0\) and \(2e_1\). Their integer
coefficients in the same declared lattice presentations are \(2\) and \(1\).
Both messages are nonzero, both exact components remain zero, and the public
observation is again shared. No function of that observation can return the
correct coefficient for both presentations. Each message remains the unique
CVP answer in its actual lattice.

\Needspace{8\baselineskip}
\begin{corollary}[Strict separation for hidden lattice recovery]
  \label{cor:lattice-member-separation}
  Fix a correct CVP oracle for the presented rank-one lattices of the
  constructed family. When the operative lattice is supplied, the oracle
  recovers its unique closest vector, equal to the operative harmonic target.
  From the shared public class and cochain alone, no solver with that same
  fixed oracle can recover the operative target for both hidden members.
  Every larger compatible family retaining those two members inherits this
  exact-recovery obstruction.
\end{corollary}

\begin{proof}
  For each supplied lattice, oracle correctness and
  \cref{prop:cellular-axis-ambiguity} identify its answer with the harmonic
  message. The two hidden members have the same public observation and
  distinct operative targets. A solver receives the same observation and
  the same oracle in either case, so its single output cannot equal both
  targets. The argument allows arbitrary fixed auxiliary information,
  including all candidate CVP answers, and imposes no running-time bound.
  Any compatible family containing the pair retains this contradiction.
\end{proof}

Solving the candidate CVP instances does not select the operative HCP answer.

\section*{Conclusion and Further Work}
\addcontentsline{toc}{section}{Conclusion and Further Work}

An explicitly presented harmonic lattice gives a precise discrete-CVP comparison, while
the constructed common-input examples prevent uniform exact recovery from
the specified public data.


The comparison suggests studying encoded HCP families with explicit
distributions on hidden presentations and Hodge-structured noise. Their
computational difficulty must be assessed within that specified model.

Discrete Morse reduction for cellular sheaf cohomology
\cite{curryGhristNanda2015morse} suggests a computational refinement.
A chain-homotopy equivalence alone does not identify the original harmonic
representative or preserve its distance. An application to the original
recovery target needs explicit retained lifts and a metric comparison.

Further work includes extending the homological model to full first-order
fibrational semantics, realizing richer families topologically, and studying
transported observations.

\appendix

\section{Paper--Lean Cross-Reference}

The shared manuscript ownership, notation, and issue register is maintained
in \lean{SEMANTIC-PARITY.md}. The companion coding paper's detailed object
and theorem ledger is \lean{papers/proximity-prize/SEMANTIC_PARITY.md}.
The appendices below record the HCP paper's own declaration and file coverage.

\begin{leanledger}

\subsection{Definitions and Interfaces}

\begin{description}[leftmargin=2.4em,style=nextline]
  \item[Matrix-level public projection and fibers.]
    Sections 2--3; \lean{PublicProjection},
    \lean{HiddenOperatorFiber}, and
    \lean{HiddenOperatorClass} in \lean{HCP.lean}.
  \item[Matrix harmonicity and Hodge instances.]
    Section 2; \lean{Hodge_Instance} and
    \lean{MatrixIsHarmonic} in \lean{HCP.lean}.
  \item[Kernel ambiguity and lossy public data.]
    Sections 2--3; \lean{SameKernel},
    \lean{KernelAmbiguous},
    \lean{PublicCompatibleHodgeAmbiguity}, and
    \lean{LossyPublicProjection} in \lean{HCP.lean}.
  \item[Hodge cipher and operative witness.]
    Section 2; \lean{Hcrypt.HCP.hodgeDecompositionEquiv},
    \lean{Hcrypt.HCP.encrypt}, \lean{Hcrypt.HCP.decrypt},
    \lean{Hcrypt.HCPBridge.Observation},
    \lean{Hcrypt.HCPBridge.HiddenHodgeWitness},
    \lean{Hcrypt.HCPBridge.operativeTarget}, and
    \lean{Hcrypt.HCPBridge.IsOperativeSolution}.
  \item[Representative-indexed public HCP solvers.]
    Section 3; \lean{Hcrypt.HCPBridge.Observation},
    \lean{Hcrypt.HCPBridge.Solver},
    \lean{Hcrypt.HCPBridge.SolverSoundOn}, and
    \lean{Hcrypt.HCPBridge.SolverExactOnHarmonic}. The direct target contract
    is \lean{Hcrypt.HCPBridge.SolverExactOnOperativeTarget}.
  \item[Explicit-lattice CVP comparison.]
    Section 3; \lean{CVP.Instance}, \lean{CVP.Solution}, and
    \lean{CVP.OracleCorrect}.
  \item[Harmonic-lattice bridge.]
    Section 3; \lean{HarmonicCVPBridge} and
    \lean{HarmonicCVPBridge.projectedTarget}.
  \item[Closest-harmonic solution predicates.]
    Section 3; \lean{ClosestHarmonicSolution} and
    \lean{ProjectionClosestHarmonicSolution}.
  \item[Hodge projection model.]
    Sections 2--3; \lean{Laplacian.harmonicProjection},
    \lean{Laplacian.residual}, and \lean{Laplacian.IsHarmonic}.
  \item[Categorical fibration copy surface.]
    Section 4; \lean{IsCartesianMorphism},
    \lean{GrothendieckFibration},
    \lean{CategoricallyEquivalentFibrations}, and
    \lean{FibrationCopy}.
  \item[Typed Helfer/HCP bridge.]
    Sections 2 and 4; \lean{Hcrypt.Helfer.Signature},
    \lean{Hcrypt.Helfer.Interpretation},
    \lean{Hcrypt.Helfer.HomotopyHomomorphism},
    \lean{Hcrypt.Helfer.HomotopyEquivalent},
    \lean{Hcrypt.Helfer.ClosedFormulaInvariance}, and
    \lean{Hcrypt.HCPBridge.Realization}.
  \item[Semantic/public equality surface.]
    Section 4; \lean{Axioms.HelferFOL},
    \lean{Axioms.PublicClassifierFactorization},
    \lean{Axioms.PublicClassifierFactorization.publicFiber_eq_iUnion_baseFibers},
    \lean{Axioms.HomotopyClass}, and \lean{Axioms.Fibration}.
  \item[Heyting/oracle public judgment.]
    Section 4; \lean{Axioms.Heyting3}, \lean{Axioms.OracleQuery},
    \lean{Axioms.OracleAnswer}, and \lean{Axioms.OracleEta2}.
  \item[Concrete finite cellular sheaf and degree-one Hodge operator.]
    Sections 2 and 4; \lean{Sheaf.OrientedTwoComplex},
    \lean{Sheaf.CellularSheaf},
    \lean{Sheaf.CellularSheaf.delta0},
    \lean{Sheaf.CellularSheaf.delta1}, and
    \lean{Sheaf.CellularSheaf.degreeOneLaplacian}.
  \item[Isometric cellular-cochain transport.]
    Section 2; \lean{Hcrypt.Sheaf.CellularSheaf.IsometricCochainEquiv},
    \lean{Hcrypt.Sheaf.CellularSheaf.IsometricCochainEquiv.degreeOneLaplacian_naturality},
    and
    \lean{Hcrypt.Sheaf.CellularSheaf.IsometricCochainEquiv.map_ker_degreeOneLaplacian}.
  \item[Snapshot-bound exact one-cochain.]
    Sections 4--5; \lean{Sheaf.CellularSheaf.SheafStateSnapshot},
    \lean{Sheaf.CellularSheaf.CochainPayloadInterpretation}, and
    \lean{Sheaf.CellularSheaf.AcceptedCommittedOneCochain}.
  \item[Integral sheaf and harmonic message lattice.]
    Sections 2--3; \lean{Sheaf.CellularSheaf.StalkLattice},
    \lean{Sheaf.CellularSheaf.IntegralStructure}, and
    \lean{Sheaf.CellularSheaf.IntegralStructure.harmonicOneLattice}.
    The restriction-preserving integral structure is supplied explicitly.
  \item[Structured Hodge sample and presented harmonic sublattice.]
    Sections 2--3; \lean{Sheaf.CellularSheaf.HodgeNoise},
    \lean{Sheaf.CellularSheaf.HodgeEncryptedCochain}, and
    \lean{Sheaf.CellularSheaf.HarmonicLatticePresentation}.
    A presented full-span sublattice is additional data, not an automatic
    identification of a real harmonic subspace with a discrete lattice.
  \item[Cochain commitment payload.]
    Section 5; \lean{Axioms.CochainCommitmentData},
    \lean{CochainCommitmentData.WellFormed}, and
    \lean{CochainCommitmentData.AddCompatible}.
  \item[Commitment security interface.]
    Section 5; \lean{Axioms.CochainCommitmentAxioms}.
  \item[Proof-carrying public artifact.]
    Section 5; \lean{TCB.HashedCommittedCochain} and
    \lean{TCB.ProofCarryingCommittedCochain}.
  \item[Bound full-sample/exact-payload observation.]
    Sections 2 and 5; \lean{TCB.HCPStackBinding},
    \lean{TCB.BoundHCPObservation}, and
    \lean{TCB.BoundHCPObservation.toObservation}.
    This witness-side record retains the message/opening and a supplied
    context-indexed binding relation; it is not itself a public wire format.
  \item[Projection chain.]
    Section 6; \lean{ProjectionChain.Chain} and
    \lean{ProjectionChain.Chain.toEvaluationProjection}.
  \item[Proximity-stage certificates.]
    Section 6; \lean{Proximity.Gates.SourceMetricGate},
    \lean{Proximity.Gates.HardGateCertificate},
    \lean{Proximity.Gates.ProjectionNaturalityGate}, and
    \lean{Proximity.Gates.NaturalityGateCertificate}. The separately defined
    \lean{Proximity.Gates.SyndromeAdequacyGate} is optional and is not a hard
    gate of the accepted-commitment projection chain.
  \item[Canonical RS/MCA endpoint.]
    Section 6; \lean{CanonicalRS.BridgeData}, \lean{CanonicalRS.MCA.Line},
    and \lean{CanonicalRS.MCA.badChallengeFraction}.
\end{description}

\subsection{Theorems and Status}

\begin{description}[leftmargin=2.4em,style=nextline]
  \item[Kernel ambiguity implies lossy public projection.]
    \lean{kernelAmbiguous_implies_lossy};
    Lean-proved.
  \item[Public-compatible Hodge ambiguity gives kernel ambiguity.]
    \lean{kernelAmbiguous_of_publicCompatibleHodgeAmbiguity};
    Lean-proved.
  \item[Helfer equivalence reaches the HCP public classifier.]
    \lean{Hcrypt.HCPBridge.samePublicClass_of_homotopyEquivalent} and
    \lean{Hcrypt.HCPBridge.publiclyEquivalent_laplacians_of_homotopyEquivalent};
    Lean-proved from the explicit closed-formula invariance and classifier
    factorization hypotheses.
  \item[Typed Hodge-extension separation gives kernel ambiguity.]
    \lean{Hcrypt.HCPBridge.publicCompatibleHodgeAmbiguity_of_kernelSeparation},
    \lean{Hcrypt.HCPBridge.kernelAmbiguous_of_kernelSeparation}, and
    \lean{Hcrypt.HCPBridge.lossyPublicProjection_of_kernelSeparation};
    Lean-proved from an explicit HCP extension-separation witness.
  \item[Hodge encryption and operative target.]
    \lean{Hcrypt.HCP.split_encrypt},
    \lean{Hcrypt.HCP.encrypt_split},
    \lean{Hcrypt.HCP.decrypt_encrypt},
    \lean{Hcrypt.HCP.decrypt_coe_eq_harmonicProjection}, and
    \lean{Hcrypt.HCPBridge.operativeTarget_is_harmonic};
    Lean-proved. The bundled bijection is
    \(\mathcal H\times\mathcal H^\perp\simeq C^1\); the paper separately
    proves the exact/coexact refinement.
  \item[Isometric transport and fixed-coordinate separation.]
    \lean{Hcrypt.Sheaf.CellularSheaf.IsometricCochainEquiv.degreeOneLaplacian_naturality},
    \lean{Hcrypt.Sheaf.CellularSheaf.IsometricCochainEquiv.isHarmonicOneCochain_iff},
    \lean{Hcrypt.HCPBridge.IsometricPresentationEquivalence.harmonic_iff_transport},
    \lean{Hcrypt.HCPBridge.IsometricPresentationEquivalence.map_matrixHarmonicSubspace},
    \lean{Hcrypt.HCPBridge.IsometricPresentationEquivalence.operativeHarmonicComponent_naturality},
    and
    \lean{Hcrypt.HCPBridge.kernelSeparationWitness_of_isometricNonstabilization};
    Lean-proved at the concrete sheaf and abstract HCP interfaces.
  \item[Concrete Hodge ambiguity exists in nontrivial dimension.]
    \lean{Hodge_Ambiguity},
    \lean{constantPublicProjection_publicCompatibleHodgeAmbiguity},
    \lean{constantPublicProjection_kernelAmbiguous}, and
    \lean{constantPublicProjection_lossy};
    Lean-proved.
  \item[Every field-linear image can be a linear kernel.]
    \lean{exists_matrix_ker_eq_range};
    Lean-proved linear-algebra construction.
  \item[Kernel realization for presented generated subspaces.]
    \lean{Hodge_GeneratedSubspace_Reduction} and
    \lean{Hodge_GeneratedSubspace_Reduction_Proof};
    Lean-proved surface.
  \item[Hodge decomposition into harmonic, exact, and coexact parts.]
    \lean{Hodge_Decomposition_Theorem};
    Lean-proved existence and pairwise orthogonality. The paper derives
    uniqueness from those orthogonal summands.
  \item[Ambiguous input contradicts public-data-only Hodge solving.]
    \lean{public_solver_eq_input_of_harmonic} and
    \lean{no_public_solver_on_ambiguous_input};
    Lean-proved lemmas.
  \item[Known-lattice CVP does not select the hidden HCP representative.]
    \lean{Hcrypt.HCPBridge.operativeTargets_ne_of_kernelSeparation},
    \lean{Hcrypt.HCPBridge.no_solver_exactOnOperativeTarget_of_kernelSeparation},
    and \lean{Hcrypt.HCPBridge.cvpOracle_insufficient_for_operativeTarget};
    Lean-proved for the direct operative-target contract. The soundness plus
    harmonic-fixed-point formulation is
    \lean{Hcrypt.HCPBridge.no_solver_on_kernelSeparation} and
    \lean{Hcrypt.HCPBridge.cvpOracle_insufficient_of_kernelSeparation}.
    The older matrix-wide theorems
    \lean{no_uniform_public_hodge_solver} and
    \lean{CVP_Oracle_Insufficient_For_Public_Hodge_Solving} remain
    generic sanity results, not the concrete HCP bridge.
  \item[Projection-model closest-harmonic surface.]
    \lean{matrix_harmonic_iff_projection_isHarmonic},
    \lean{Hodge_Projection_Is_ClosestHarmonic}, and
    \lean{harmonicProjection_is_closestHarmonic};
    Lean-proved.
  \item[Discrete CVP minimizers are preserved by harmonic projection.]
    \lean{HarmonicCVPBridge.squaredDistance_projection_decomposition},
    \lean{HarmonicCVPBridge.isClosestLatticeVector_iff_projectedTarget},
    and \lean{HarmonicCVPBridge.solution_iff_projectedTarget};
    Lean-proved.
  \item[Fibration copies preserve the categorical fibration surface.]
    \lean{copy_categoricallyEquivalent_thm} and
    \lean{FibrationCopy_is_Grothendieck};
    Lean-proved categorical interface.
  \item[Exact sheaf gates supply \(d^{k+1}d^k=0\) and Laplacian-from-coboundary.]
    \lean{Gates.structural_preconditions_of_sheaf_exact_gates};
    Lean-proved alias/interface.
  \item[Concrete cellular coboundaries form a complex.]
    \lean{Sheaf.CellularSheaf.delta1_comp_delta0} and
    \lean{Sheaf.CellularSheaf.isOneCocycle_of_isExactOneCochain};
    Lean-proved from the oriented cell incidences and restriction-composition
    law.
  \item[Concrete degree-one Hodge characterization and orthogonality.]
    \lean{Sheaf.CellularSheaf.degreeOneLaplacian_energy},
    \lean{Sheaf.CellularSheaf.ker_degreeOneLaplacian},
    \lean{Sheaf.CellularSheaf.exactOneSubspace_isOrtho_harmonicOneSubspace},
    \lean{Sheaf.CellularSheaf.harmonicOneSubspace_isOrtho_coexactOneSubspace},
    and \lean{Sheaf.CellularSheaf.exactOneSubspace_isOrtho_coexactOneSubspace};
    Lean-proved.  The separately defined
    \lean{Sheaf.CellularSheaf.kernelNormalizedDegreeOneLaplacian} is proved to
    have the same kernel, not asserted to equal the true Hodge operator or to
    have its spectrum.
  \item[Exact payloads determine openings modulo \(\ker\delta_0\).]
    \lean{Sheaf.CellularSheaf.openingEquivalent_iff_sub_mem_ker},
    \lean{Sheaf.CellularSheaf.decoded_isExactOneCochain}, and
    \lean{Sheaf.CellularSheaf.decoded_isOneCocycle};
    Lean-proved for the concrete cellular-sheaf commitment object.
  \item[Cellular coboundaries preserve the supplied integral structure.]
    \lean{Sheaf.CellularSheaf.IntegralStructure.delta0_mem_integralOneCochains}
    and
    \lean{Sheaf.CellularSheaf.IntegralStructure.delta1_mem_integralTwoCochains};
    Lean-proved from the stalk-lattice restriction-preservation fields.
  \item[Structured samples recover their integral harmonic message.]
    \lean{Sheaf.CellularSheaf.HodgeNoise.value_mem_harmonic_orthogonal} and
    \lean{Sheaf.CellularSheaf.HodgeEncryptedCochain.harmonicOneProjection_observed};
    Lean-proved for the supplied integral sheaf and exact/coexact opening.
    The ambient observed sample is not asserted to be integral or exact.
  \item[Presented harmonic-lattice samples have the encoded point as unique CVP solution.]
    \lean{Sheaf.CellularSheaf.HarmonicLatticePresentation.squaredDistance_observed_decomposition},
    \lean{Sheaf.CellularSheaf.HarmonicLatticePresentation.solution_cvpInstance_iff_eq_point},
    and
    \lean{Sheaf.CellularSheaf.HarmonicLatticePresentation.existsUnique_solution_cvpInstance};
    conditional on an explicit harmonic-lattice presentation and a sample
    constructed from one of its lattice points. This does not supply a
    general CVP algorithm or derive the presentation from public data.
  \item[An actual integral family instantiates the semantic HCP ambiguity.]
    \lean{HCP.AxisFamily.shared_observation_distinct_nonzero_targets},
    \lean{HCP.AxisFamily.no_exact_public_solver},
    \lean{HCP.AxisFamily.cvp_oracle_insufficient}, and
    \lean{HCP.AxisFamily.cvp_solution_iff_message};
    the two cellular sheaves, integral structures, actual Hodge matrices,
    common structured-noise observation and harmonic lattices are constructed.
    Homological-test invariance is proved in the actual bounded-below real
    cochain model. The semantics is a stated fragment, and the oracle
    corollary retains its oracle-existence premise. The observation is
    \((p,b)\), without a full storage-view or security-game claim.
    \lean{weighted_public_observations_eq} and
    \lean{no_exact_encoded_message_decoder} in the same namespace additionally
    separate the explicit integer coefficients \(2\) and \(1\) at one
    shared observation, rather than only separate ambient coordinates.
  \item[The actual bound observation agrees with the semantic operative target.]
    \lean{TCB.HCPRealization.BoundObservationAlignment.operativeTarget_eq_message}
    and \lean{cvp_solution_iff_operativeTarget} in the same namespace;
    derived from actual-operator action and classifier alignment, retaining
    the original sample, payload, snapshot and artifact. A selected richer
    public view requires its own observation-factorization proof.
  \item[The bound observation retains payload/snapshot identity and harmonic recovery.]
    \lean{TCB.BoundHCPObservation.artifact_wellFormed},
    \lean{TCB.BoundHCPObservation.payload_boundToSnapshot}, and
    \lean{TCB.BoundHCPObservation.harmonicProjection_cochainObject};
    Lean-proved from the record's accepted exact component, artifact binding,
    decoded-component equality, and supplied observation-binding relation.
    No modular serializer or concrete public semantic realization follows
    from constructing this conditional record.
  \item[Well-formed payloads have nonempty edge data and supported modulus.]
    \lean{CochainCommitmentData.edge_values_nonempty_of_wellFormed} and
    \lean{CochainCommitmentData.modulus_gt_one_of_wellFormed};
    Lean-proved.
  \item[Proof-carrying artifacts expose the committed payload.]
    \lean{TCB.proofCarrying_cochain_wellFormed};
    TCB boundary theorem.
  \item[Projection chain composes to one evaluation projection.]
    \lean{ProjectionChain.Chain.toEvaluationProjection};
    Lean-proved plumbing.
  \item[Syndrome adequacy turns zero local syndrome into codeword/global candidate.]
    \lean{Proximity.Gates.SyndromeAdequacyGate.codeword_of_wellFormed_zeroSyndrome};
    stage-local interface.
  \item[Stage naturality composes without rechecking.]
    \lean{Proximity.Gates.ProjectionNaturalityGate.of_stageCertificate};
    Lean-proved composition of the cochain/plane and plane/RS restriction and
    syndrome squares.
  \item[Final proximity evidence yields calibrated end-to-end readiness.]
    \lean{Proximity.Gates.CalibratedFinalCertificate.toCalibratedEndToEndReady};
    Lean-proved plumbing.
  \item[Canonical RS list/MCA bounds enter at a cited boundary.]
    \lean{CanonicalRS.CitedBounds.bounds};
    cited external theorem package.
  \item[Canonical MCA bad-challenge fraction is bounded after projection.]
    \lean{CanonicalRS.MCA.projectionBadChallengeFraction_le};
    Lean-proved from cited bounds.
\end{description}

\end{leanledger}

\section{Lean Source Coverage by File}

This appendix is the file-by-file coverage audit for the translation. The body
of the paper gives the mathematical definitions and proofs; this appendix
records where each Lean declaration family enters the written account.
Import-only aggregate files are listed with the files they collect.

\begin{leanledger}

\subsection{Build and Aggregate Files}

\begin{description}[leftmargin=2.4em,style=nextline]
  \item[\lean{lakefile.lean}.]
    Build configuration; no mathematical declaration is translated as a
    theorem.
  \item[\lean{Sheaf.lean}, \lean{Axioms.lean}, \lean{Gates.lean}, \lean{TCB.lean},
    \lean{Consensus.lean}, \lean{Proximity.lean}, and
    \lean{Proximity/CanonicalRS.lean}.]
    Aggregate import surfaces. Their mathematical content is exactly the
    union of the component files listed below.
\end{description}

\subsection{HCP, CVP, and Fibrational Algebra}

\begin{description}[leftmargin=2.4em,style=nextline]
  \item[\lean{HCP.lean}.]
    Section 2 contains the matrix-level HCP vocabulary, public fibers,
    ambiguity hierarchy, typed Helfer realization, representative-bound
    observations and witnesses, the Hodge-coordinate cipher, and the unique
    operative recovery target. It also records isometric presentation
    transport, harmonic-projection naturality, fixed-coordinate
    nonstabilization, classifier nonfaithfulness, and the resulting
    no-exact-representative-decoder theorem. The direct public recovery
    contract is
    \lean{Hcrypt.HCPBridge.SolverExactOnOperativeTarget}.
  \item[\lean{Hardness.lean}.]
    Sections 2--4 contain consequences and comparisons for the HCP interface
    imported from \lean{HCP.lean}. The concrete ambiguity theorems
    \lean{Hodge_Ambiguity}, \lean{Hodge_Ambiguity_Proof},
    \lean{constantPublicProjection_publicCompatibleHodgeAmbiguity},
    \lean{constantPublicProjection_kernelAmbiguous}, and
    \lean{constantPublicProjection_lossy} supply the matrix sanity example
    discussed after \cref{prop:cvp-insufficient}, rather than a concrete typed
    HCP separation witness. The presented-linear-image comparison
    declarations \lean{Hodge_Solver_Correct},
    \lean{exists_matrix_ker_eq_range},
    \lean{Hodge_GeneratedSubspace_Reduction}, and
    \lean{Hodge_GeneratedSubspace_Reduction_Proof} are
    \cref{prop:hcp-to-cvp}. The decomposition theorem
    \lean{Hodge_Decomposition_Theorem} is
    \cref{prop:hodge-decomposition}. The public-solver predicates
    \lean{PublicHodgeInstance}, \lean{PublicHodgeSolver},
    \lean{Hcrypt.HCPBridge.operativeTargets_ne_of_kernelSeparation},
    \lean{Hcrypt.HCPBridge.no_solver_exactOnOperativeTarget_of_kernelSeparation},
    \lean{Hcrypt.HCPBridge.cvpOracle_insufficient_for_operativeTarget},
    \lean{Hcrypt.HCPBridge.no_solver_on_kernelSeparation}, and
    \lean{Hcrypt.HCPBridge.cvpOracle_insufficient_of_kernelSeparation} are
    \cref{prop:cvp-insufficient}. The closest-harmonic predicates
    \lean{ClosestHarmonicSolution}, \lean{ProjectionClosestHarmonicSolution},
    \lean{matrix_harmonic_iff_projection_isHarmonic},
    \lean{Hodge_Projection_Is_ClosestHarmonic}, and
    \lean{harmonicProjection_is_closestHarmonic} record the continuous
    projection surface. The discrete bridge declarations
    \lean{HarmonicCVPBridge},
    \lean{HarmonicCVPBridge.squaredDistance_projection_decomposition},
    \lean{HarmonicCVPBridge.isClosestLatticeVector_iff_projectedTarget}, and
    \lean{HarmonicCVPBridge.solution_iff_projectedTarget} are
    \cref{prop:hodge-projection-is-cvp}. The categorical declarations
    \lean{IsCartesianMorphism}, \lean{GrothendieckFibration},
    \lean{CategoricallyEquivalentFibrations}, \lean{Wrapper}, \lean{Fibration},
    \lean{FibrationCopy}, \lean{ToWrapper}, \lean{FromWrapper},
    \lean{copy_categoricallyEquivalent_thm}, and \lean{FibrationCopy_is_Grothendieck}
    are \cref{prop:fibration-copy-equivalence}.
  \item[\lean{CVP.lean}.]
    Section 3 translates \lean{AmbientVector}, \lean{ExplicitLattice},
    \lean{latticeVector},
    \lean{InExplicitLattice}, \lean{IsClosestLatticeVector},
    \lean{squaredEuclideanDistance}, \lean{Instance}, \lean{Solution},
    \lean{Oracle}, \lean{LatticeOracleCorrect}, \lean{OracleCorrect},
    \lean{SolverCorrect}, \lean{solverCorrect_iff_oracleCorrect},
    \lean{HardnessAssumption}, and
    \lean{zero_lattice_oracle_correct} as the known-lattice comparison
    surface.
  \item[\lean{Helfer/Signature.lean},
    \lean{Helfer/HomotopyEquivalence.lean},
    \lean{Helfer/Invariance.lean}, \lean{Axioms/HelferFOL.lean},
    \lean{Axioms/Fibration.lean}, and
    \lean{Axioms/HomotopyClass.lean}.]
    Section 4 formalizes multi-sorted signatures, finite-product
    interpretations, coherent homotopy homomorphisms, and componentwise
    homotopy equivalence for cofibrant-fibrant interpretations. It also
    extracts the Helfer-derived path/diagonal and fibrational interface, then
    adds the project-specific public classifiers and public fibers, classifier
    factorization through a base projection, class-membership witnesses,
    transport, semantic well-formedness clauses, fiber systems, ephemeral
    fiber specifications, and nonfaithful classifiers. This is not a complete
    formalization of Helfer's model-categorical construction or
    formula-in-context invariance proof; the closed-formula consequence of
    Theorem~12.3 is an explicit Lean hypothesis. The local equivalence, symmetry, transitivity,
    membership, and nonfaithfulness lemmas justify the public fiber language
    used in HCP.
\end{description}

\subsection{Foundational, Sheaf, Spectral, and Judgment Axioms}

\begin{description}[leftmargin=2.4em,style=nextline]
  \item[\lean{Sheaf/CellComplex.lean}, \lean{Sheaf/Cellular.lean},
    \lean{Sheaf/Hodge.lean}, \lean{Sheaf/Equivalence.lean}, and
    \lean{Sheaf/Commitment.lean}.]
    Sections 2, 4, and 5 translate the concrete finite oriented two-complex;
    finite real coordinate stalks and face-to-coface maps; the typed spaces
    \(C^0,C^1,C^2\); concrete \(\delta_0\) and \(\delta_1\); the proof that
    \(\delta_1\delta_0=0\); degree-zero global sections, degree-one cocycles,
    and exact degree-one cochains; the true degree-one Hodge operator
    \(\delta_0\delta_0^*+\delta_1^*\delta_1\); its energy, kernel, and
    orthogonality theorems; the distinct kernel-normalized operator;
    degreewise isometric cochain equivalences and their adjoint, Laplacian,
    harmonic-subspace, and kernel transport theorems; immutable
    sheaf-state snapshots; snapshot-bound payload interpretation; accepted
    exact one-cochains; and the quotient of openings by \(\ker\delta_0\).
  \item[\lean{Sheaf/Integral.lean}.]
    Sections 2--3 and 5 use \lean{StalkLattice}, its additive carrier and
    \lean{Preserves} predicate, and \lean{IntegralStructure}. The latter
    selects integral stalk lattices preserved by the cellular restrictions.
    The induced \lean{integralZeroCochains}, \lean{integralOneCochains}, and
    \lean{integralTwoCochains} are additive subgroups;
    \lean{harmonicOneLattice} is the integral degree-one subgroup intersected
    with the harmonic subspace. The two integral-coboundary preservation
    theorems are proved for this supplied structure. No canonical modular
    encoding/decoding is constructed here.
  \item[\lean{Sheaf/Distance.lean} and \lean{Sheaf/CochainSemantics.lean}.]
    Distance to the actual exact subspace is nonnegative, vanishes exactly
    on exact cochains, and is stable under Euclidean perturbations. The
    semantics adapter intertwines the actual cochain operators and decoder,
    identifies snapshot binding, and derives exactness, cocycle and acceptance
    transport. Other legacy ambient gates remain explicit. The carrier
    equivalence is not asserted to preserve norms.
  \item[\lean{Sheaf/HarmonicLattice.lean}.]
    Sections 2--3 use \lean{HodgeNoise}, \lean{HodgeEncryptedCochain},
    \lean{observed}, and harmonic projection recovery for
    \(b=h+\delta_0a+\delta_1^*c\).
    \lean{HarmonicLatticePresentation} additionally supplies an explicit
    full-span harmonic sublattice, membership proofs and its real-span
    equality. Its \lean{point}, \lean{encrypt}, and \lean{cvpInstance}
    constructions lead to the squared-distance decomposition and unique
    closest-point result for those encoded samples. The presentation need
    not contain every integral harmonic cochain. Noise sampling, a secrecy
    game and modular wire realization remain separate obligations.
  \item[\lean{Sheaf/ObservationQuotient.lean}.]
    The full kernel-normalized residual equals the Hodge-noise value.
    \Cref{prop:exact-component-obstruction} is formalized by the coexact-ray,
    residual-nonfactorization and far-preservation obstruction theorems.
    Their source carrier is unrestricted Hodge noise, and the extraction
    must factor through its decoded exact component; no raw-byte
    factorization or unrestricted production-profile claim is inferred.
  \item[\lean{Axioms/Foundation.lean} and \lean{Gates/Exact.lean}.]
    Section 4 translates \lean{Foundation} and
    \lean{FoundationWellFormed}: the cochain-complex premise, coboundary
    linearity, Hodge--Laplacian formation, Hodge decomposition, and Hodge
    orthogonality.  Commitment authentication/hiding, topology connectivity,
    spectral policy, and key-image/freshness conditions are separate
    interfaces rather than fields of \lean{Foundation}. The exact-gate
    aliases \lean{ExactGateBundle}, \lean{SheafExactGateBundle},
    \lean{FoundationWithSheafExactGates}, and
    \lean{structural_preconditions_of_sheaf_exact_gates} are
    \cref{prop:exact-structural-preconditions}.
  \item[\lean{Axioms/SheafCoherence.lean}.]
    Section 4 translates restriction-shape compatibility, edge naturality,
    triangle-face commutativity, sheaf cochain complexes, Hodge coherence,
    finite simplicial nerves, \lean{SheafCoherence},
    \lean{SheafCoherenceWellFormed},
    typed degree-zero through degree-two coboundaries,
    \lean{CommittedCochainWellFormed},
    \lean{RepresentsExactOneCochain}, \lean{IsOneCocycle},
    \lean{oneCocycle_of_committedCochainWellFormed},
    \lean{FoundationWithSheafCoherence}, and
    \lean{coherence_implies_structural_preconditions}.
  \item[\lean{Axioms/Grassmannian.lean}.]
    Sections 4 and 6 translate the Grassmannian metric and geometry objects:
    \lean{GrassmannianDistance}, \lean{chordal_leq_geodesic},
    \lean{KarcherMean}, \lean{SheafLaplacian}, \lean{CechNerve},
    \lean{PluckerEligibility}, \lean{ConsensusGeometry}, and
    \lean{ConsensusGeometryWellFormed}.
  \item[\lean{Axioms/SpectralBounds.lean} and \lean{Gates/Spectral.lean}.]
    Section 4 translates the constants
    \(\sqrt2/2\) and \(2\sqrt2\), the predicates
    \lean{SpectralGapInPoVBounds}, \lean{NearPoVSpectralBoundary},
    \lean{TsirelsonPredicate}, \lean{NearTsirelsonBoundary}, the aliases
    \lean{SpectralGapGate} and \lean{NearSpectralBoundary}, the geometry
    witness \lean{ByConstructionPoVGeometry}, and the theorems
    \lean{byConstruction_gap_in_bounds} and
    \lean{byConstruction_passes_spectral_gate}.
  \item[\lean{Tsirelson.lean}.]
    Section 4 translates the Grassmannian/Laplacian family
    \lean{Grassmannian_Laplacians}, \leanOmega{_G},
    \lean{LaplacianSpace}, \lean{SpectralGap}, \leanOmega{_1},
    \lean{projector_map}, \lean{Classical_Strategy},
    \lean{Quantum_Strategy}, and \lean{MaxEigenvalue}; the infinitude and
    injectivity results \lean{R_infinite}, \lean{infinite_lines},
    \lean{Grassmannian_Laplacian_injective},
    \lean{projector_injective}, \lean{laplacian_injective},
    \leanOmega{_1_infinite}, \lean{R_uncountable},
    \leanOmega{_1_infinite_v2}, and \leanOmega{_1_Nonempty}; the Laplacian-space
    inclusions \lean{laplacian_properties},
    \lean{laplacian_is_hermitian},
    \leanOmega{_1_subset_LaplacianSpace},
    \leanOmega{_1_subset_LaplacianSpace_v3},
    \leanOmega{_1_subset_LaplacianSpace_v5},
    \leanOmega{_1_subset_LaplacianSpace_explicit_v2}, and
    \leanOmega{_1_subset_LaplacianSpace_explicit}; the classical/quantum
    statement forms \lean{Classical_Bound_Statement} and
    \lean{Quantum_Bound_Statement}; and the empty-family sanity checks
    \lean{Classical_Bound_empty_family} and
    \lean{Quantum_Bound_empty_family}.
  \item[\lean{Axioms/Heyting.lean} and
    \lean{Axioms/OracleModality.lean}.]
    Section 4 translates the three-valued Heyting algebra, meet, join,
    implication, negation, failure of double-negation elimination, abstract
    Heyting collapse, nonfaithful public truth, spectral judgments, consensus
    actions, oracle queries and answers, \(\Omega_{\mathrm{nv}}\), witness
    collapses, quotient equivalence, sheafification, oracle forcing,
    \(\eta_1/\eta_2\) oracle modalities, Lawvere--Tierney idempotence,
    boundary judgment behavior, canonical consensus admissibility, oracle
    pipelines, and the oracle/consensus bridge theorems.
  \Needspace{6\baselineskip}
  \item[\lean{Axioms/StageAdmissibility.lean},
    \lean{Gates/Stage.lean}, and \lean{Gates/Oracle.lean}.]
    This appendix records the stage and oracle aliases: pipeline stages, stage
    admissibility, transaction-critical invariant bundles, Python-stage
    snapshots, refinement predicates, and the theorem that a refined
    transaction-critical pass is admissible.
\end{description}

\subsection{Laplacian Geometry and Search Spaces}

\begin{description}[leftmargin=2.4em,style=nextline]
  \item[\lean{Laplacian/Hodge.lean} and \lean{Laplacian.lean}.]
    Sections 2--3 translate harmonic projection, residual vector, residual
    norm, harmonicity, projection membership, nonnegativity, and all
    equivalences between zero residual, membership in the harmonic subspace,
    and harmonicity.
  \item[\lean{Laplacian/Projection.lean}.]
    Section 3 translates the projection Laplacian
    \(\Delta_W=I-P_W\) and its pointwise application formula.
  \item[\lean{Laplacian/Spectral.lean}.]
    Section 4 translates sheaf Laplacians, spectral fingerprints, spectral
    equivalence, scalar Laplacians, kernel and eigenvalue lemmas, scalar
    fingerprints, the legacy relation defined by equal kernel dimension, and
    the theorem that spectrum is not determined by nullity. That relation
    does not assert Helfer homotopy equivalence.
  \item[\lean{Laplacian/Likelihood.lean}.]
    Section 4 records classical selector and density-based likelihood
    interfaces, Gaussian likelihoods, minimum spectral gaps, integrated
    likelihood scores and Gaussian factorization lemmas. Legacy names such
    as ``quantum adversary'' and \lean{success_prob_*} supply neither a
    quantum computation model nor HCP recovery probabilities. The corrected
    Euclidean integrated scores are not identified with the legacy
    supremum-norm integrals. These interfaces are outside the selected entry.
  \item[\lean{Laplacian/Grassmannian.lean},
    \lean{Laplacian/Grassmannian/Injectivity.lean},
    \lean{Laplacian/Grassmannian/SearchSpace.lean}, and
    \lean{Laplacian/Grassmannian/ParametricFamily.lean}.]
    Sections 4 and 6 translate harmonic subspaces, the map to projection
    Laplacians, the families \(\Omega_r\) and \(\Omega_1\), membership in
    Laplacian space, nonemptiness of \(\Omega_k\), injectivity of the
    Grassmannian Laplacian map, and the parametric construction of families of
    \(k\)-planes.
\end{description}

\subsection{Commitments, Public Artifacts, and TCB Boundary}

\begin{description}[leftmargin=2.4em,style=nextline]
  \item[\lean{Axioms/CochainCommitment/Payload.lean},
    \lean{Axioms/CochainCommitment/Blinded.lean},
    \lean{Axioms/CochainCommitment/Security.lean}, and
    \lean{Axioms/CochainCommitment.lean}.]
    Section 5 translates the canonical modulus, byte lengths, wire length,
    \lean{CochainCommitmentData}, edge-count matching, nonempty edge data,
    supported modulus, well-formedness, commitment sums, common-topology and
    additive compatibility, zero and addition operations, payload lemmas,
    exact-payload authentication, wrapper hiding, context binding,
    additivity/security interfaces, the typed exact-payload to
    Hodge-blinded-value bridge, exact--harmonic orthogonality and decomposition
    theorems, blinded and vector commitment axioms, cohomological MAC axioms,
    and the canonical commitment construction from foundational geometry plus
    separately supplied commitment-specific hypotheses.
  \item[\lean{TCB/ProofCarrying.lean}, \lean{TCB/Runtime.lean},
    \lean{TCB/SNARK.lean}, \lean{TCB/CompressionBoundary.lean},
    \lean{TCB/ObservationBoundary.lean}, and
    \lean{TCB/WitnessVerifier.lean}.]
    Section 5 and this appendix record hash words, hashed committed cochains,
    proof-carrying cochains, runtime specification boundaries, runtime
    acceptance predicates, SNARK-carrying cochains, compression contracts,
    compressed public data, observation boundaries, gap-decision certificates,
    witness-compression envelopes, public state updates, and theorems stating
    that accepted artifacts expose well-formed cochain and oracle predicates.
  \item[\lean{Helfer/CochainModel.lean}, \lean{Sheaf/HCPAxisFamily.lean},
    \lean{HCP/AxisFamily.lean}.]
    Section 3's concrete family. The first module constructs the real cochain
    Quillen model, retains the actual cellular differentials and proves
    homological-test invariance for all typed interpretations. The second
    constructs integral sheaves, samples, operator matrices, chain isometry,
    spectrum and harmonic lattices. The third derives the semantic realization,
    common observation, distinct nonzero targets and both positive and negative
    CVP comparisons from those same objects. No full first-order semantics,
    cryptographic sampling experiment or runtime profile is constructed.
  \item[\lean{TCB/HCPRealization.lean}, \lean{TCB/HCPAxisFamily.lean}.]
    Sections 3 and 5's application join. Actual matrix action gives the
    cellular kernel/projector equality; snapshot-classifier agreement gives
    the original observation's hidden witness and operative message. The
    existing CVP comparison is transported to that target. The concrete family
    instantiates the operator alignment and can retain an original artifact
    under supplied acceptance, decoding and envelope bindings. The actual
    computed positive gap also satisfies the spectral gate. Complete public
    storage-view equality is not inferred from the class equality.
  \item[\lean{TCB/HCPObservation.lean}.]
    Sections 2 and 5 use \lean{HCPStackBinding} and
    \lean{BoundHCPObservation} to join the full Hodge sample to its accepted
    exact committed component, artifact leaf digest, immutable snapshot root
    and lifecycle context. The \lean{latticeObject}, \lean{cochainObject},
    \lean{hashObject}, \lean{snapshotRoot}, and \lean{toObservation} accessors
    retain those same objects. Artifact well-formedness, snapshot binding and
    harmonic recovery are proved from the explicit record fields. The
    classifier, decoder and observation-binding relation are supplied;
    constructing the record does not instantiate their public realization.
    The record contains private witness data and is not a serialization or
    public disclosure prescription.
  \item[\lean{Axioms/ZKP.lean}.]
    Section 5 records the zero-knowledge proof object and its well-formedness
    predicate as an assumed proof-carrying artifact interface.
\end{description}

\subsection{Projection Chain, Proximity Gates, and RS/MCA Endpoint}

The universal standalone coding results are owned by the companion Paper 3
and its detailed source ledger. Its completed generic distance, real-shell,
phase-split and all-parameter family constructions reside in
\lean{Proximity/Core/Distance.lean},
\lean{Proximity/Core/RadiusShell.lean},
\lean{Proximity/Capacity/Realization.lean},
\lean{Proximity/Capacity/ListPhase.lean},
\lean{Proximity/ListDecoding/PhaseSplit.lean},
\lean{Proximity/MCA/RealThreshold.lean}, and
\lean{Proximity/Application/DistanceCalibration.lean}.
This appendix records their HCP-specific consumers without repeating the
standalone coding proofs.

\begin{description}[leftmargin=2.4em,style=nextline]
  \item[\lean{Proximity/Problem.lean}, \lean{Proximity/Public.lean},
    \lean{Proximity/Projection.lean}, \lean{Proximity/PluckerShadow.lean},
    \lean{Proximity/GrassmannianPlucker.lean}, and
    \lean{Proximity/PluckerCochain/PublicBinding.lean}.]
    Section 6 records the problem-facing public projection vocabulary, the
    abstract Pl\"ucker coordinate-system and public-shadow contracts, and the
    interfaces that prepare the concrete Pl\"ucker coordinate and finite-code
    datum.
  \item[\lean{Proximity/Application/GapBoundary.lean}.]
    Section 6.1 formalizes disjointness of the closed promises under a strict
    gap, the equal-radius boundary obstruction, and strict-gap views retaining
    the existing carrier, score and radii. The classical specification solver
    characterizes logical consistency; it supplies neither an executable or
    efficient solver nor a concrete application profile.
  \item[\lean{Proximity/Grassmannian.lean},
    \lean{Proximity/GrassmannianPlucker.lean}, and
    \lean{Proximity/GrassmannianPlucker/Concrete.lean}.]
    Section 6 translates harmonic Grassmannian points, exterior coordinates,
    frame witnesses, decomposability, framed systems, raw and normalized
    Pl\"ucker coordinates, soundness, projective equivalence, projective
    normalization, projective framed systems, and projective sound inputs.
  \item[\lean{Proximity/Concrete/GrassmannianPlucker/Reconstruction.lean}
    and \lean{TCB/PluckerObservation.lean}.]
    Section 6 proves both concrete exterior-algebra soundness inputs,
    reconstructs the operative subspace and harmonic projector from a genuine
    full projective point, and proves full-shadow equality exactly when the
    ranked points agree. Identification with \(\ker\Delta_x\) gives the
    operative harmonic target; digest inversion and efficient real-coordinate
    algorithms are outside these results.
  \item[\lean{Proximity/CochainBridge.lean},
    \lean{Proximity/PluckerCochain.lean},
    \lean{Proximity/PluckerCochain/Concrete.lean}, and
    \lean{Proximity/PluckerCochain/PublicBinding.lean}.]
    Section 6 translates valid frames, Pl\"ucker-to-cochain instances, local
    syndrome targets, codeword predicates, evaluation projections, invariant
    checklists, payload encoders, commitment-sum residuals, payload
    realizations, hashed Pl\"ucker coordinates, observed frame bindings, and the
    theorems that valid Pl\"ucker coordinates bind to sheaf-coherent exact
    one-cochains, hence to accepted committed cocycles and local codewords.
  \item[\lean{Proximity/CochainPlane.lean},
    \lean{Proximity/CochainPlane/Core.lean},
    \lean{Proximity/CochainPlane/Concrete.lean},
    \lean{Proximity/CochainPlane/ReedSolomon.lean}, and
    \lean{Proximity/CochainReedSolomon.lean}.]
    Section 6 translates plane projections, plane extraction contexts,
    local nondegeneracy, alignment, projected planes, plane targets,
    plane-to-evaluation projections, typed restriction and syndrome naturality
    at the cochain/plane and plane/RS transitions,
    local-codeword-to-target-codeword theorems, separate plane and RS metric
    handoffs, faithful-on-region plane encodings, interpretations, and RS
    received words.
  \item[\lean{Proximity/ProjectionChain.lean} and
    \lean{Proximity/ProjectionChain/Core.lean}.]
    Section 6 translates the complete projection chain, including accessors
    for committed cochain, plane projection, received word, and the composite
    evaluation projection.
  \item[\lean{Proximity/Gates.lean}, \lean{Proximity/Gates/Stage.lean},
    \lean{Proximity/Gates/Final.lean}, and
    \lean{Proximity/GateEvidence.lean}.]
    Section 6 translates only proximity gates: source metric, Pl\"ucker metric,
    semantic alignment, optional syndrome adequacy, projective-coordinate
    respect, three-stage noncollapse through cochain and plane radii,
    far-region separation, runtime binding, structural bundles, hard-gate and
    derived naturality certificates, calibrated final certificates, nested
    witness-local state buses, stage checklists, stage certificates, radius
    calibrations, final readiness, evidence slots, evidence databases, and the
    binding/completeness lemmas for proximity evidence.
  \item[\lean{Proximity/RS/ListDecoding.lean},
    \lean{Proximity/RS/ReedSolomon.lean},
    \lean{Proximity/RS/CorrelatedAgreement.lean},
    \lean{Proximity/RS/Target.lean}, and \lean{Proximity/RS.lean}.]
    Section 6 translates words, block codes, Hamming and relative Hamming
    balls, list-size bounds, interleaved and additive codes, RS parameters,
    vector codewords, correlated-agreement and mutual-correlated-agreement
    predicates, bad-parameter sets, bridge conditions, calibrated
    \(\HE\)-to-RS bridges, and pullback-correctness theorems.
  \item[\lean{Proximity/CanonicalRS/Parameters.lean},
    \lean{Proximity/CanonicalRS/PowerDomain.lean},
    \lean{Proximity/CanonicalRS/Encoding.lean},
    \lean{Proximity/CanonicalRS/Metric.lean},
    \lean{Proximity/CanonicalRS/Target.lean},
    \lean{Proximity/CanonicalRS/FarSeparation.lean},
    \lean{Proximity/CanonicalRS/BridgeData.lean},
    \lean{Proximity/CanonicalRS/BridgeWitness.lean},
    \lean{Proximity/CanonicalRS/BridgeCertificate.lean},
    \lean{Proximity/CanonicalRS/CitedBounds.lean},
    \lean{Proximity/CanonicalRS/MCA.lean},
    \lean{Proximity/CanonicalRS/MCA/Line.lean}, and
    \lean{Proximity/CanonicalRS/MCA/Concrete.lean}.]
    Section 6 translates the canonical finite field, length, degree bound,
    rate, Johnson and capacity radii, smooth domains, power-root domain,
    canonical encoding, the proof that the unbounded integer encoding is not
    globally injective, Hamming-to-code metrics, target bounds, far-separation
    witnesses, bridge data and theorem slots, bridge certificates, cited
    BCHKS-style bounds, MCA block code, challenge and bad-challenge fractions,
    and the bound on the projected bad-challenge fraction.
\end{description}

\newpage
\subsection{Subsequent Protocol and Implementation Boundary}

The bounded state/transition, candidate-indexed evidence and typed storage
modules under \lean{Protocol/Ledger/}, \lean{Protocol/Evidence/}, and
\lean{Protocol/Storage/} belong to Paper 2;
their completion status is recorded in \lean{papers/paper2/SCOPE.md}.
They are downstream consumers of the mathematical objects above.

\begin{description}[leftmargin=2.4em,style=nextline]
  \item[\lean{Axioms/Consensus.lean}, \lean{Axioms/Ledger.lean},
    \lean{Axioms/Coregulator.lean}, and
    \lean{Axioms/PythonHomotopyRefinement.lean}.]
    This appendix records consensus, ledger, coregulator, and Python homotopy
    refinement as boundary objects. Their full protocol semantics belong to
    the subsequent protocol paper, but their predicates are named here so the
    Paper 1 mathematical objects have explicit downstream interfaces.
  \item[\lean{Consensus/Geometry.lean},
    \lean{Consensus/Reliability.lean}, and
    \lean{Consensus/ProofOfValue.lean}.]
    This appendix records Pl\"ucker eligibility and consensus geometry aliases,
    reliability evidence and bounded reliability scores, proximity quality,
    Proof-of-Value gates, Proof-of-Value inputs, scores with and without
    reliability/proximity, and the score bounds. These are protocol-facing
    uses of the Section 4 and Section 6 mathematical predicates.
  \item[\lean{Chain.lean}.]
    This appendix records validators, transactions, blocks, protocol states, spectral
    bounds, topology and freshness predicates, block validity, tips, heights,
    extension, transaction and block acceptance, ledger-state updates,
    key-image preservation theorems, validator eligibility, cumulative gap,
    preference, transaction composition and identities, quadratic weight,
    Sybil-resistance weighting, chain-growth monotonicity, fork-resolution
    transitivity, and connectedness of eligible validators. The detailed
    protocol architecture remains in the subsequent protocol paper.
\end{description}

\end{leanledger}


\clearpage
\bibliographystyle{alpha}
\bibliography{references}

\end{document}
